\chapter{The Geometry of Numbers}
\begin{center}\itshape
A convex region of the plane, symmetric about $0$, with area greater than $4$, always contains a point of $\Z^2$\\
other than $0$. This innocent statement proves that every class group is finite.
\end{center}

In 1881 the Paris Academy set a prize question: find the number of ways to write an integer as a sum of five squares.
The prize went in April 1883 to two entrants. One was Henry Smith of Oxford, the leading English number theorist, who in fact
had published the answer years before, and who died two months before the award. The other was Hermann Minkowski, a
student in K\"onigsberg, who was seventeen when he sent in his memoir. The commission could not separate the two, and
gave the prize to both; in England, where Smith was far better known, the decision caused an outcry. Camille Jordan, of
the commission, wrote to the young man: ``Travaillez, Monsieur, \`a devenir un g\'eom\`etre \'eminent'' (work, Sir, to
become an eminent geometer).

He did, and the word \emph{geometer} turned out to be prophetic. Thinking about quadratic forms, Minkowski realised that
many questions about integers are questions about \emph{lattice points}, the points with integer coordinates, and that one
can answer them by drawing. His key observation, published in 1891 and expanded in the book \emph{Geometrie der Zahlen}
(1896), is the theorem at the top of this page and its $n$-dimensional version: a convex body that is symmetric and large
compared with a lattice must contain a lattice point other than $0$. It sounds like a puzzle for a rainy afternoon.

Volume~1 proved that the class group of a quadratic field is finite with a pigeonhole argument (Volume~1, Lemma~11.12 and
Theorem~11.13), and the bound it gave was poor. Minkowski's idea is to look at $\mathcal O_K$ as a lattice in $\R^n$
(Chapter~3, Figure~\ref{fig:mink}) and use volume: choose a body on which every point has small norm, and you get an
element of small norm in any ideal you like. From this one trick come four theorems. The class group is finite, with a
bound small enough to compute class numbers by hand. Every number field other than $\Q$ has a ramified prime. There are only
finitely many number fields of bounded discriminant (Hermite). And, in Chapter~6, the unit group is as large as it can be
(Dirichlet). The story has a modern sequel: finding short vectors in lattices of high dimension is so hard that it now
protects internet traffic against quantum computers.

\begin{reminder}[from Volume 6 and Chapters 3--4]
\emph{Lebesgue measure} $\lambda$ on $\R^n$ (Volume~6, Chapter~6) is countably additive and invariant under translation, and
gives the box $\prod[a_i,b_i]$ its volume $\prod(b_i-a_i)$. Closed sets and open sets are measurable. \emph{Linear maps}
(Volume~6, Chapter~8): $\lambda(T(A))=|\det T|\,\lambda(A)$ for $T$ linear. \emph{Fubini--Tonelli and polar coordinates}
(Volume~6, Chapter~8): volumes can be computed slice by slice, and for $g\ge0$ measurable,
$\int_{\R^2}g(|z|)\dd z=2\pi\int_0^\infty g(\rho)\rho\dd\rho$. \emph{Chapter~3.} The signature $(r_1,r_2)$, $n=r_1+2r_2$;
the discriminant $d_K$ of an integral basis. \emph{Chapter~4.} Unique factorisation of ideals, the norm
$N(I)=[\mathcal O:I]$, multiplicative, with $N(\alpha\mathcal O)=|N(\alpha)|$; the class group $\operatorname{Cl}(K)$;
a prime ramifies if and only if it divides $d_K$.
\end{reminder}

\section{Lattices and Minkowski's theorem}

\begin{definition}
A \emph{lattice} in $\R^n$ is a subgroup $L=A\Z^n$, where $A\in\mathrm{GL}_n(\R)$; the columns of $A$ are a
\emph{basis} of $L$. Its \emph{fundamental parallelotope} is $F=A[0,1)^n$, and its \emph{covolume} is
$\operatorname{covol}(L)=\lambda(F)=|\det A|$.
\end{definition}

The covolume does not depend on the basis: two bases differ by a matrix of $\mathrm{GL}_n(\Z)$, of determinant $\pm1$. If
$L'\subseteq L$ are lattices, $[L:L']$ is finite and $\operatorname{covol}(L')=[L:L']\operatorname{covol}(L)$: in a basis
of $L$, $L'$ is $C\Z^n$ with $C\in M_n(\Z)$, $[L:L']=|\det C|$ (Volume~2, Corollary~7.8), and $\det(AC)=\det A\det C$.

Two facts about a lattice $L=A\Z^n$ are used all the time. It is \emph{discrete}: a bounded set $B$ contains finitely
many points of $L$, because $A^{-1}B$ is bounded and contains finitely many points of $\Z^n$. And the translates $F+\ell$,
$\ell\in L$, \emph{tile} $\R^n$: every $x$ is uniquely $f+\ell$ with $f\in F$, $\ell\in L$ (write $A^{-1}x=y$ and split
each coordinate of $y$ into its integer and fractional parts).

\begin{lemma}[Blichfeldt]\label{lem:blichfeldt}
Let $L$ be a lattice and $S\subseteq\R^n$ measurable with $\lambda(S)>\operatorname{covol}(L)$. Then there are $x\ne y$ in
$S$ with $x-y\in L$.
\end{lemma}
\begin{proof}
By the tiling, $S$ is the disjoint union of the countably many sets $S_\ell=S\cap(F+\ell)$, $\ell\in L$. The translates
$S_\ell-\ell$ lie in $F$, and $\sum_\ell\lambda(S_\ell-\ell)=\sum_\ell\lambda(S_\ell)=\lambda(S)>\lambda(F)$ by translation
invariance and countable additivity. If the sets $S_\ell-\ell$ were pairwise disjoint, their union would be a subset of
$F$ of measure $>\lambda(F)$. So two of them meet: $x-\ell=y-\ell'$ with $x\in S_\ell$, $y\in S_{\ell'}$, $\ell\ne\ell'$.
Then $x\ne y$ (they lie in different tiles) and $x-y=\ell-\ell'\in L$.
\end{proof}

\begin{theorem}[Minkowski's convex body theorem]\label{thm:minkowski}
Let $L\subseteq\R^n$ be a lattice and $S\subseteq\R^n$ a measurable set that is convex and symmetric ($x\in S\Rightarrow
-x\in S$).
\begin{enumerate}[label=(\alph*)]
\item If $\lambda(S)>2^n\operatorname{covol}(L)$, then $S$ contains a point of $L\setminus\{0\}$.
\item If $S$ is moreover compact, the same holds when $\lambda(S)\ge2^n\operatorname{covol}(L)$.
\end{enumerate}
\end{theorem}
\begin{proof}
(a) The set $\frac12S$ has measure $2^{-n}\lambda(S)>\operatorname{covol}(L)$. By Blichfeldt there are $x\ne y$ in
$\frac12S$ with $x-y\in L$. Now $2x$ and $-2y$ are in $S$ (symmetry), so their midpoint $x-y$ is in $S$ (convexity), and
it is a non-zero point of $L$.

(b) For each $k\ge1$, $(1+\frac1k)S$ is convex, symmetric, of measure $(1+\frac1k)^n\lambda(S)>2^n\operatorname{covol}(L)$,
so it contains some $\ell_k\in L\setminus\{0\}$. All the $\ell_k$ lie in the bounded set $2S$, which contains finitely
many lattice points, so some $\ell\ne0$ equals $\ell_k$ for infinitely many $k$. Then $\ell\in(1+\frac1k)S$ for
arbitrarily large $k$, that is $\frac{k}{k+1}\ell\in S$, and as $S$ is closed, $\ell\in S$.
\end{proof}

\begin{figure}[ht]
\centering
\fig{ch03-minkowski}
\caption{A lattice $L$ with its fundamental parallelogram (shaded), and a symmetric convex body $S$ of area
$4.4\operatorname{covol}(L)>2^2\operatorname{covol}(L)$. Minkowski's theorem guarantees the rust points. Shrink $S$ to
half its size, and the proof's Blichfeldt step finds two points of $\frac12S$ differing by a lattice vector.}
\label{fig:mink3}
\end{figure}

\begin{think}
Each hypothesis is needed. Find a symmetric but non-convex set, and a convex but non-symmetric one, of area as large
as you like, containing no point of $\Z^2\setminus\{0\}$. Then find a compact convex symmetric set of area exactly $4$
whose only lattice points are on its boundary.
\end{think}

The constant $2^n$ is sharp: the open cube $(-1,1)^n$ has volume $2^n$ and contains no non-zero point of $\Z^n$.

\begin{example}[Fermat's two squares, again]\label{exa:twosq}
Let $p\equiv1\pmod 4$ and $u$ with $u^2\equiv-1\pmod p$ (Volume~1, Chapter~5). The set $L=\{(a,b)\in\Z^2:a\equiv ub\pmod p\}$
is a subgroup of index $p$ in $\Z^2$ (it is the kernel of the onto map $(a,b)\mapsto a-ub\bmod p$), hence a lattice of
covolume $p$. The disc $a^2+b^2<2p$ has area $2\pi p>4p$, so it contains $(a,b)\in L\setminus\{0\}$. Then $a^2+b^2\equiv
u^2b^2+b^2\equiv0\pmod p$ and $0<a^2+b^2<2p$, so $a^2+b^2=p$. Exercise~\ref{ex:foursq} gives four squares the same way.
\end{example}

\section{The Minkowski embedding}

Let $K$ have signature $(r_1,r_2)$, with real embeddings $\sigma_1,\dots,\sigma_{r_1}$ and one embedding $\tau_j$ chosen
in each pair of complex ones. The \emph{Minkowski embedding} is
\[
j:K\to\R^{r_1}\times\C^{r_2}\cong\R^n,\qquad x\mapsto\bigl(\sigma_1(x),\dots,\sigma_{r_1}(x),\tau_1(x),\dots,
\tau_{r_2}(x)\bigr),
\]
where $\C$ is identified with $\R^2$ by $z\mapsto(\operatorname{Re}z,\operatorname{Im}z)$. It is an injective ring
homomorphism (coordinatewise operations on the right), and $\Q$-linear.

\begin{proposition}\label{prop:covol}
$j(\mathcal O_K)$ is a lattice of covolume $2^{-r_2}\sqrt{|d_K|}$, and for an ideal $I$, $j(I)$ is a lattice of covolume
$2^{-r_2}\sqrt{|d_K|}\,N(I)$.
\end{proposition}
\begin{proof}
Let $\omega_1,\dots,\omega_n$ be an integral basis and $A$ the real $n\times n$ matrix whose columns are the
$j(\omega_k)$. Its rows are $\sigma_i(\omega_k)$ for the real embeddings, and the pairs $\operatorname{Re}\tau_j(\omega_k)
=\frac{\tau_j+\bar\tau_j}2(\omega_k)$, $\operatorname{Im}\tau_j(\omega_k)=\frac{\tau_j-\bar\tau_j}{2i}(\omega_k)$. Replacing
each such pair of rows by the rows $\tau_j(\omega_k)$, $\bar\tau_j(\omega_k)$ is a row operation with determinant
$\det\begin{psmallmatrix}1/2&1/2\\1/(2i)&-1/(2i)\end{psmallmatrix}=\frac{-1}{2i}$, of absolute value $\frac12$. So
$|\det A|=2^{-r_2}|\det(\sigma(\omega_k))|=2^{-r_2}\sqrt{|d_K|}$, which is not zero: $j(\mathcal O_K)=A\Z^n$ is a
lattice. Finally $j(I)$ is a subgroup of index $[\mathcal O:I]=N(I)$ in $j(\mathcal O_K)$, a lattice of covolume
$N(I)$ times larger.
\end{proof}

For $\Z[\sqrt{-5}]$, $n=2$, $r_2=1$, $d_K=-20$, the lattice $\Z+\Z\sqrt5\,i\subseteq\C$ has covolume $\sqrt5=\frac12\sqrt{20}$,
and the ideal $P=(2,1+\sqrt{-5})$, of norm $2$, has covolume $2\sqrt5$ (Figure~\ref{fig:ideal3}).

\begin{figure}[ht]
\centering
\fig{ch03-ideal}
\caption{$\Z[\sqrt{-5}]$ in $\C$ (all dots) and its ideal $P=(2,1+\sqrt{-5})$ (rust dots: $a+b\sqrt{-5}$ with $a\equiv b
\pmod 2$). The fundamental parallelogram of $P$, spanned by $2$ and $1+\sqrt{-5}$, has twice the area of that of $\Z[\sqrt{-5}]$,
spanned by $1$ and $\sqrt{-5}$: the covolume grows by the index $N(P)=2$.}
\label{fig:ideal3}
\end{figure}

\section{The Minkowski bound and the class group}

To find elements of small norm, we need a convex symmetric set on which the norm is small. The norm of $x$ is
$|N(x)|=\prod_i|\sigma_i(x)|\prod_j|\tau_j(x)|^2$, a product of $n$ absolute values when complex coordinates are counted
twice, and the arithmetic--geometric mean inequality bounds a product by a sum.

\begin{lemma}\label{lem:volume}
For $t>0$ let $B_t=\{(y,z)\in\R^{r_1}\times\C^{r_2}:\sum_i|y_i|+2\sum_j|z_j|\le t\}$. It is compact, convex and
symmetric, and
\[
\lambda(B_t)=2^{r_1}\Bigl(\frac\pi2\Bigr)^{r_2}\frac{t^n}{n!}.
\]
For $x\in K$ with $j(x)\in B_t$, $|N(x)|\le t^n/n^n$.
\end{lemma}
\begin{proof}
$B_t$ is closed and bounded, symmetric, and convex because $(y,z)\mapsto\sum|y_i|+2\sum|z_j|$ is a norm. Write
$V_{r_1,r_2}(t)$ for its volume. Scaling by $t$ multiplies volume by $t^n$, so $V_{r_1,r_2}(t)=c_{r_1,r_2}t^n$, and
$c_{1,0}=2$ (the interval $[-1,1]$), $c_{0,1}=\pi/4$ (the disc of radius $\frac12$). Slicing along one real coordinate
(Fubini),
\[
V_{r_1+1,r_2}(t)=\int_{-t}^tV_{r_1,r_2}(t-|y|)\dd y=2c_{r_1,r_2}\int_0^t(t-y)^n\dd y=\frac{2c_{r_1,r_2}}{n+1}t^{n+1}.
\]
Slicing along one complex coordinate and using polar coordinates, then $u=2\rho$,
\begin{align*}
V_{r_1,r_2+1}(t)&=\int_{2|z|\le t}c_{r_1,r_2}(t-2|z|)^n\dd z\\
&=2\pi c_{r_1,r_2}\int_0^{t/2}(t-2\rho)^n\rho\dd\rho\\
&=\frac\pi2c_{r_1,r_2}\int_0^t(t-u)^nu\dd u\\
&=\frac\pi2\cdot\frac{c_{r_1,r_2}\,t^{n+2}}{(n+1)(n+2)},
\end{align*}
the last integral by one integration by parts. Starting from $c_{0,0}=1$ (with $V_{0,0}=1$ in dimension $0$), each real
coordinate multiplies by $2$ and each complex one by $\frac\pi2$, while the factorials accumulate to $\frac1{n!}$.

For the norm, apply the arithmetic--geometric mean inequality to the $n$ numbers $|\sigma_i(x)|$ and $|\tau_j(x)|$ (each
$|\tau_j(x)|$ twice): $|N(x)|^{1/n}\le\frac1n\bigl(\sum|\sigma_i(x)|+2\sum|\tau_j(x)|\bigr)\le\frac tn$.
\end{proof}

\begin{theorem}[Minkowski's bound]\label{thm:minkbound}
Let
\[
M_K=\Bigl(\frac4\pi\Bigr)^{r_2}\frac{n!}{n^n}\sqrt{|d_K|}.
\]
Every ideal $I$ contains a non-zero $\alpha$ with $|N(\alpha)|\le M_K\,N(I)$, and every ideal class contains an ideal of
norm at most $M_K$. In particular the class group is finite.
\end{theorem}
\begin{proof}
Choose $t$ with $\lambda(B_t)=2^n\operatorname{covol}(j(I))$, that is, by Proposition~\ref{prop:covol} and
Lemma~\ref{lem:volume},
\[
2^{r_1}\Bigl(\frac\pi2\Bigr)^{r_2}\frac{t^n}{n!}=2^{n-r_2}\sqrt{|d_K|}\,N(I),
\]
so
\[
t^n=2^{n-r_1-2r_2}\Bigl(\frac4\pi\Bigr)^{r_2}n!\sqrt{|d_K|}\,N(I)=\Bigl(\frac4\pi\Bigr)^{r_2}n!\sqrt{|d_K|}\,N(I).
\]
By Theorem~\ref{thm:minkowski}(b) there is $\alpha\in I$, $\alpha\ne0$, with $j(\alpha)\in B_t$, and then $|N(\alpha)|
\le t^n/n^n=M_KN(I)$.

Now let $C$ be an ideal class and $I$ an ideal in the inverse class $C^{-1}$ (every class contains ideals: multiply a
fractional ideal by a suitable integer). Take $\alpha$ as above. Since $\alpha\mathcal O\subseteq I$, $\alpha\mathcal O=IJ$
for an ideal $J$ (Corollary~\ref{cor:divides}), whose class is $C$, and $N(J)=|N(\alpha)|/N(I)\le M_K$. Finally there are
only finitely many ideals of norm at most $M_K$: each contains its norm (Proposition~\ref{prop:normideal}), and an
integer $m$ lies in only finitely many ideals (Proposition~\ref{prop:noeth}).
\end{proof}

The theorem turns the class group into a finite computation: $\operatorname{Cl}(K)$ is generated by the classes of the
primes of norm at most $M_K$, which lie above the prime numbers $p\le M_K$, and Theorem~\ref{thm:kd} finds them. Then one
looks for relations among them, coming from elements of small norm.

\begin{example}[{$\Q(\sqrt{-5})$}]
$n=2$, $r_2=1$, $d_K=-20$: $M_K=\frac4\pi\cdot\frac24\sqrt{20}=\frac{4\sqrt5}\pi\approx2.85$. So the class group is
generated by the primes above $2$: $2\mathcal O=P^2$ with $P=(2,1+\sqrt{-5})$, not principal since $a^2+5b^2=2$ has no
solution. As $P^2=(2)$ is principal, $\operatorname{Cl}(K)\cong\Z/2\Z$ and $h_K=2$, as Volume~1 found by a longer route.
\end{example}

\begin{example}[{$\Q(\sqrt[3]2)$ and $\Q(\sqrt[3]{10})$}]
For $\Q(\sqrt[3]2)$, $(r_1,r_2)=(1,1)$, $d_K=-108$, $M_K=\frac4\pi\cdot\frac6{27}\sqrt{108}\approx2.94$, and the only
prime of norm $\le2$ is $(\sqrt[3]2)$, principal: $h_K=1$. For $\Q(\sqrt[3]{10})$, $d_K=-300$ gives $M_K\approx4.90$, so we
must look at the primes above $2$ and $3$; Exercise~\ref{ex:h10} finds that they are all principal, so $h_K=1$ again.
For $\Q(\sqrt[3]7)$, on the other hand, $h_K=3$.
\end{example}

\section{Discriminants: Minkowski and Hermite}

\begin{theorem}[Minkowski]\label{thm:dkbig}
If $K\ne\Q$, then
\[
\sqrt{|d_K|}\ \ge\ \Bigl(\frac\pi4\Bigr)^{r_2}\frac{n^n}{n!}\ \ge\ \Bigl(\frac\pi4\Bigr)^{n/2}\frac{n^n}{n!}>1.
\]
In particular $|d_K|>1$, so at least one prime ramifies in $K$.
\end{theorem}
\begin{proof}
Every ideal has norm $\ge1$, so applying Theorem~\ref{thm:minkbound} to $I=\mathcal O$ gives $1\le|N(\alpha)|\le M_K$,
which is the first inequality. The second uses $\frac\pi4<1$ and $r_2\le\frac n2$. For the third let
$a_n=(\frac\pi4)^{n/2}\frac{n^n}{n!}$. Then $a_2=\frac\pi4\cdot2=\frac\pi2>1$ and $\frac{a_{n+1}}{a_n}=\sqrt{\frac\pi4}
\bigl(1+\frac1n\bigr)^n\ge\sqrt{\frac\pi4}\cdot2>1$. So $a_n>1$ for $n\ge2$. A prime dividing $d_K$ ramifies by
Theorem~\ref{thm:ramdisc}.
\end{proof}

Since $\frac{n^n}{n!}\ge\frac{e^n}{e\sqrt n}$ roughly (Stirling, Volume~4), $|d_K|$ grows at least exponentially in the
degree: a field of large degree has a large discriminant. Exercise~\ref{ex:discbound} gives the clean bound
$|d_K|\ge\frac\pi3\bigl(\frac{3\pi}4\bigr)^{n-1}$. The ramification statement is the arithmetic analogue of a
topological fact: $\Q$ has no unramified extension, just as the sphere has no connected unramified (covering) map onto it
other than the identity (Chapter~13 takes up this analogy).

\begin{theorem}[Hermite]\label{thm:hermite}
For every $X>0$ there are only finitely many number fields $K\subseteq\C$ with $|d_K|\le X$.
\end{theorem}
\begin{proof}
By Theorem~\ref{thm:dkbig}, $(\frac\pi4)^{n/2}\frac{n^n}{n!}\le\sqrt X$, and the left side tends to infinity with $n$, so
the degree is bounded. Fix $n$ and the signature $(r_1,r_2)$. We find $\theta\in\mathcal O_K$ generating $K$ whose
conjugates are bounded in terms of $X$ only; its minimal polynomial then has bounded integer coefficients, and there are
finitely many such polynomials, each with finitely many roots.

\emph{If $r_1\ge1$}, let $S$ be the set of $(y,z)$ with $|y_1|\le C$, $|y_i|\le\frac12$ for $2\le i\le r_1$, and
$|z_j|\le\frac12$ for all $j$. It is compact, convex, symmetric, of volume $2C\cdot(\frac\pi4)^{r_2}$. Choose $C$ so that
this equals $2^{n-r_2}\sqrt{|d_K|}$; then $C\le C_X=2^{n-r_2-1}(\frac4\pi)^{r_2}\sqrt X$. Theorem~\ref{thm:minkowski}
gives $\theta\in\mathcal O_K\setminus\{0\}$ with $j(\theta)\in S$. As $|N(\theta)|\ge1$ and all conjugates but
$\sigma_1(\theta)$ have absolute value $\le\frac12$, $|\sigma_1(\theta)|\ge2^{n-1}>\frac12$. So $\sigma_1(\theta)$ is
different from every other conjugate $\sigma(\theta)$.

\emph{If $r_1=0$}, let $S$ be the set where $|\operatorname{Re}z_1|\le\frac12$, $|\operatorname{Im}z_1|\le C$, and $|z_j|
\le\frac12$ for $j\ge2$, of volume $2C(\frac\pi4)^{r_2-1}$, and choose $C$ as before. We get $\theta$ with $|\tau_1(\theta)|^2
\ge4^{r_2-1}\ge1$; as $|\operatorname{Re}\tau_1(\theta)|\le\frac12$, $\tau_1(\theta)$ is not real, so it differs from
its conjugate $\bar\tau_1(\theta)$, and it differs from the other conjugates, which have absolute value $\le\frac12$.

In both cases some embedding takes on $\theta$ a value taken by no other embedding. Now the values $\sigma_i(\theta)$
are the roots of $\chi_\theta=m_\theta^{n/d}$, $d=[\Q(\theta):\Q]$ (Proposition~\ref{prop:charpoly}), so each of them is
taken $n/d$ times. Hence $n/d=1$ and $K=\Q(\theta)$. All conjugates of $\theta$
have absolute value at most $C_X+1$ (the complex ones satisfy $|\tau_1(\theta)|^2\le C_X^2+\frac14$), so the coefficients of
$m_\theta$, elementary symmetric functions of the conjugates, are integers bounded by $\binom nk(C_X+1)^k$.
\end{proof}

Hermite's theorem raises a counting question: how many fields of degree $n$ have $|d_K|\le X$? For $n=2$ the answer is
about $\frac6{\pi^2}X$ (squarefree numbers, Volume~5), for $n=3$ about $\frac X{3\zeta(3)}$ (Davenport--Heilbronn), and
Malle's conjecture predicts the answer for every Galois group. Volume~14 returns to these \emph{arithmetic statistics}.

\section{Exercises}

\begin{exercise}[$\star$]\label{ex:convexity}
Answer the think box: give a symmetric non-convex set and a convex non-symmetric set, both of infinite area, containing no
point of $\Z^2\setminus\{0\}$, and a compact convex symmetric set of area $4$ whose only non-zero lattice points are on its
boundary.
\end{exercise}

\begin{exercise}[$\star\star$]\label{ex:foursq}
\emph{Four squares.} Let $p$ be an odd prime. Show that there are $u,v$ with $u^2+v^2+1\equiv0\pmod p$. Show that
$L=\{(a,b,c,d)\in\Z^4:a\equiv uc+vd,\ b\equiv ud-vc\pmod p\}$ has covolume $p^2$ and that $a^2+b^2+c^2+d^2\equiv0\pmod p$
on $L$. The ball of radius $r$ in $\R^4$ has volume $\frac{\pi^2}2r^4$; deduce that $p$ is a sum of four squares.
\end{exercise}

\begin{exercise}[$\star\star$]\label{ex:h23}
Show that the class group of $\Q(\sqrt{-23})$ is cyclic of order $3$, generated by a prime above $2$.
\end{exercise}

\begin{exercise}[$\star\star$]\label{ex:h47}
Show that the class group of $\Q(\sqrt{-47})$ is cyclic of order $5$.
\end{exercise}

\begin{exercise}[$\star$]\label{ex:h1}
Show that $\Q(\sqrt{-19})$, $\Q(\sqrt{-43})$ and $\Q(\sqrt{-67})$ have class number $1$. Do the same for $\Q(\sqrt{-163})$
(the bound is about $8.1$). What does this say about the values of $n^2+n+41$ for $0\le n\le39$?
\end{exercise}

\begin{exercise}[$\star\star$]\label{ex:sqrt10}
Show that $\Q(\sqrt{10})$ has class number $2$.
\end{exercise}

\begin{exercise}[$\star\star$]\label{ex:h10}
Show that $\Q(\sqrt[3]{10})$ has class number $1$. \emph{(Use the integral basis $1,\theta,\beta$ of
Example~\ref{exa:cuberoot10}, and look for elements of norm $2$ and $3$.)}
\end{exercise}

\begin{exercise}[$\star\star$]\label{ex:discbound}
Show that $|d_K|\ge\frac\pi3\bigl(\frac{3\pi}4\bigr)^{n-1}$ for $n\ge2$. Deduce that $|d_K|\ge3$ for quadratic fields and
$|d_K|\ge13$ for cubic fields (the true minima are $3$ and $23$). Which number field of degree $n\le4$ could have
$|d_K|<100$?
\end{exercise}

\begin{exercise}[$\star\star$]\label{ex:quadmin}
\emph{A quadratic form version.} Let $Q(x,y)=ax^2+bxy+cy^2$ be positive definite with integer coefficients and
discriminant $D=b^2-4ac<0$. Show with Minkowski's theorem that $Q$ represents a non-zero value at most
$\frac2\pi\sqrt{|D|}$. Compare with the bound $\sqrt{|D|/3}$ given by Gauss's reduction (Volume~1, Chapter~7).
\end{exercise}

\begin{exercise}[$\star\star\star$]\label{ex:simult}
\emph{Simultaneous approximation (Dirichlet).} Let $\alpha_1,\dots,\alpha_k$ be real and $N\ge1$. Show that there are integers
$q$ with $1\le q\le N^k$ and $p_1,\dots,p_k$ with $|q\alpha_i-p_i|\le\frac1N$ for all $i$. \emph{(A box in $\R^{k+1}$.)}
\end{exercise}

\begin{beauty}
Minkowski's theorem is a few lines of measure theory, and yet it decides the finiteness of class groups, the existence of
ramified primes and, in the next chapter, the structure of the unit group. The reason is that $\mathcal O_K$ is not just
a ring but a \emph{lattice} with a natural geometry, and arithmetic questions become questions about short vectors. The
twentieth century took this further. Siegel computed the average number of lattice points in a body over the space of
all lattices; Hermite's and Minkowski's constants for the densest lattice packings of spheres were found in dimensions
$8$ and $24$ by Viazovska (Fields Medal, 2022), with modular forms (Chapter~17). And the hardness of finding short
vectors in high-dimensional lattices is now the basis of post-quantum cryptography.
\end{beauty}

\section{Hints}

Try each exercise for a while before reading its hint, and read the hint before the solution.

\begin{hint}{ex:convexity}
Think of a thin cross along the axes but avoiding them; a half-plane; the square $[-1,1]^2$.
\end{hint}
\begin{hint}{ex:foursq}
The sets $\{u^2\}$ and $\{-1-v^2\}$ modulo $p$ each have $\frac{p+1}2$ elements. $L$ is the kernel of a map
$\Z^4\to(\Z/p\Z)^2$ that is onto.
\end{hint}
\begin{hint}{ex:h23}
$M_K\approx3.05$. The integers are $x+y\omega$, $\omega=\frac{1+\sqrt{-23}}2$, with norm $x^2+xy+6y^2$. Look for elements
of norm $6$ and $8$.
\end{hint}
\begin{hint}{ex:h47}
$M_K\approx4.36$. The norm form is $x^2+xy+12y^2$; it takes the values $12$ and $32$ but not $2$, $3$, $6$, $8$.
\end{hint}
\begin{hint}{ex:h1}
For $d\equiv5\pmod8$, $2$ is inert. Check that the small primes are inert, with Legendre symbols.
\end{hint}
\begin{hint}{ex:sqrt10}
$M_K=\sqrt{10}$. The norm $x^2-10y^2$ is never $\pm2$ or $\pm3$ (look modulo $5$), but $N(4+\sqrt{10})=6$.
\end{hint}
\begin{hint}{ex:h10}
$M_K\approx4.90$. Kummer--Dedekind with $\theta$ works at $2$, and with $\beta$ at $3$. The norm of $\theta-a$ is
$10-a^3$, and $\beta$ has norm $3$.
\end{hint}
\begin{hint}{ex:discbound}
Use $r_2\le\frac n2$ and show that $b_n=(\frac\pi4)^n(\frac{n^n}{n!})^2$ satisfies $b_{n+1}/b_n\ge\frac{3\pi}4$.
\end{hint}
\begin{hint}{ex:quadmin}
$Q(x,y)\le r$ is an ellipse of area $\frac{2\pi r}{\sqrt{|D|}}$.
\end{hint}
\begin{hint}{ex:simult}
Apply Theorem~\ref{thm:minkowski}(b) to $\Z^{k+1}$ and the box $|x_0|\le N^k$, $|x_0\alpha_i-x_i|\le\frac1N$.
\end{hint}

\section{Solutions}

\begin{solution}{ex:convexity}
\idea Each hypothesis rules out one kind of escape: a set made of thin pieces avoiding the lattice for convexity, a set
on one side of a line for symmetry, and points on the boundary for the strict inequality.
\writeup \emph{Symmetric, not convex.} Take the set of points with $0<x<1$ or $-1<x<0$, that is $\{0<|x|<1\}$. It is symmetric, of infinite area, and its
points have $x\notin\Z$, so it contains no lattice point; it is not convex (it misses the segment joining $(\frac12,0)$
and $(-\frac12,0)$ at $x=0$).

\emph{Convex, not symmetric.} The strip $\{0<x<1\}$ is convex, of infinite area, and contains no lattice point, as $x$
would be an integer strictly between $0$ and $1$.

\emph{Area exactly $4$.} The square $[-1,1]^2$ is compact, convex, symmetric, of area $4$, and its non-zero lattice points
$(\pm1,0)$, $(0,\pm1)$, $(\pm1,\pm1)$ are on its boundary; its interior $(-1,1)^2$ has area $4$ and contains no non-zero
lattice point, so part (a) of Theorem~\ref{thm:minkowski} needs the strict inequality.
\end{solution}

\begin{solution}{ex:foursq}
\idea Same as Example~\ref{exa:twosq}: a sublattice of index $p^2$ on which the quadratic form vanishes modulo $p$, and a
ball large enough to contain a non-zero point of it, but small enough that the only multiple of $p$ available is $p$.
\writeup The squares modulo $p$ take $\frac{p+1}2$ values, so do the numbers $-1-v^2$, and $\frac{p+1}2+\frac{p+1}2>p$: two of
them coincide, $u^2\equiv-1-v^2$. The map $\Z^4\to(\Z/p\Z)^2$, $(a,b,c,d)\mapsto(a-uc-vd,\ b-ud+vc)$, is onto (take
$c=d=0$), so its kernel $L$ has index $p^2$ and covolume $p^2$. On $L$,
\[
a^2+b^2\equiv(uc+vd)^2+(ud-vc)^2=(u^2+v^2)(c^2+d^2)\equiv-(c^2+d^2)\pmod p,
\]
so $a^2+b^2+c^2+d^2\equiv0$. The open ball $a^2+b^2+c^2+d^2<2p$ has volume $\frac{\pi^2}2(2p)^2=2\pi^2p^2>16p^2=2^4
\operatorname{covol}(L)$, since $2\pi^2\approx19.7$. So it contains a non-zero point of $L$, whose sum of squares is a
multiple of $p$ in $(0,2p)$: it is $p$. By Euler's four-square identity (Volume~2, Chapter~12), products of sums of four
squares are sums of four squares, and $2=1+1+0+0$; so every positive integer is a sum of four squares.
\end{solution}

\begin{solution}{ex:h23}
\idea Below the Minkowski bound only the primes above $2$ and $3$ matter. Elements of norm $6$ and $8$ give the relations,
and the absence of elements of norm $2$, $3$, $4$ (other than $\pm2$) shows that the class of a prime above $2$ has order
exactly $3$.
\writeup Here $d_K=-23$, $r_2=1$, $M_K=\frac4\pi\cdot\frac24\sqrt{23}=\frac2\pi\sqrt{23}\approx3.05$. Since $-23\equiv1
\pmod 8$, $2$ splits: $2\mathcal O=P\bar P$ (Theorem~\ref{thm:kd} with $X^2-X+6\equiv X(X-1)\pmod2$). Since $-23\equiv1\pmod3$
is a square, $3$ splits: $3\mathcal O=Q\bar Q$ ($X^2-X+6\equiv X(X-1)\pmod 3$). These four primes generate
$\operatorname{Cl}(K)$. The norm of $x+y\omega$ is $x^2+xy+6y^2=(x+\frac y2)^2+\frac{23}4y^2$, which is $\ge\frac{23}4>5$
when $y\ne0$; so the values $2$, $3$ are not taken, and $4$ only by $\pm2$. Hence $P,Q$ are not principal, and $P^2$ is not
principal ($P^2\ne P\bar P=(2)$).

The element $\omega$ has norm $6$, so $(\omega)$ is a product of a prime above $2$ and one above $3$; renaming, $(\omega)=PQ$,
so $[Q]=[P]^{-1}$. The element $1+\omega=\frac{3+\sqrt{-23}}2$ has norm $1+1+6=8$, so $(1+\omega)$ is a product of three
primes above $2$. It is not divisible by $(2)=P\bar P$, since $\frac{1+\omega}2\notin\mathcal O$; so it is $P^3$ or $\bar
P^3$. Either way $[P]^3=1$ (as $[\bar P]=[P]^{-1}$). Since $[P]\ne1$, $[P]$ has order $3$, and $\operatorname{Cl}(K)$,
generated by $[P]$, $[\bar P]=[P]^{-1}$, $[Q]=[P]^{-1}$, $[\bar Q]=[P]$, is cyclic of order $3$.
\end{solution}

\begin{solution}{ex:h47}
\idea As for $-23$: primes above $2$ and $3$; the element of norm $32$ shows that $[P]^5=1$, and the missing norms show
that no smaller power is trivial.
\writeup $M_K=\frac2\pi\sqrt{47}\approx4.36$, so the primes above $2$ and $3$ generate. As $-47\equiv1\pmod8$ and $-47\equiv1
\pmod 3$, both split: $2=P\bar P$, $3=Q\bar Q$ (the polynomial $X^2-X+12$ has the roots $0,1$ modulo $2$ and modulo $3$).
The norm $x^2+xy+12y^2$ is $\ge\frac{47}4$ for $y\ne0$, so the values below $12$ are only the squares $x^2$: no element has
norm $2$, $3$, $6$ or $8$, and the elements of norm $4$, $9$ are $\pm2$, $\pm3$.

Now $\omega$ has norm $12$, so $(\omega)$ is a product of two primes above $2$ and one above $3$, not divisible by $(2)$
($\frac\omega2\notin\mathcal O$); renaming, $(\omega)=P^2Q$, so $[Q]=[P]^{-2}$. And $4+\omega=\frac{9+\sqrt{-47}}2$ has
norm $16+4+12=32$; it is not divisible by $2$, so $(4+\omega)=P^5$ or $\bar P^5$, and $[P]^5=1$. The class group is
generated by $[P]$ (as $[\bar P]=[P]^{-1}$, $[Q]=[P]^{-2}$, $[\bar Q]=[P]^2$), so its order divides $5$; and $[P]\ne1$
since nothing has norm $2$. So it is cyclic of order $5$.
\end{solution}

\begin{solution}{ex:h1}
\idea If every prime below the bound is inert, the only ideals of small norm are generated by integers.
\writeup For $d=-19,-43,-67,-163$, $M_K=\frac2\pi\sqrt{|d|}\approx2.77,\ 4.17,\ 5.21,\ 8.13$. All four $d$ are $\equiv5
\pmod8$, so $2$ is inert ($X^2-X+\frac{1-d}4$ has no root modulo $2$, its constant term being odd). For an odd prime
$p<M_K$, $p$ is inert when $\left(\frac dp\right)=-1$: for $d=-43$, $-43\equiv2\pmod 3$; for $d=-67$, $-67\equiv2\pmod3$ and
$-67\equiv3\pmod5$, a non-square; for $d=-163$, $-163\equiv2\pmod3$, $\equiv2\pmod5$, $\equiv5\pmod7$, all non-squares.
So every prime of norm $\le M_K$ is $(p)$, principal, and $h_K=1$.

For $d=-163$, $\omega=\frac{1+\sqrt{-163}}2$ and the norm of $x+y\omega$ is $x^2+xy+41y^2$, which is $\ge\frac{163}4>40$
when $y\ne0$. So no element has norm a prime $p<41$. Since $h_K=1$, a prime $p<41$ that split, $p\mathcal O=P\bar P$,
would have $P=(\alpha)$ with $N(\alpha)=p$: impossible. The only ramified prime is $163$. So every prime $p<41$ is inert.
Now let $0\le n\le39$ and $m=n^2+n+41=N(n+\omega)$. If $m$ were composite, it would have a prime factor $p\le\sqrt m<41$
(as $m\le39^2+39+41=1601<41^2$). An inert $p$ dividing $N(n+\omega)=(n+\omega)(n+\bar\omega)$ divides one of the two
factors, as $p\mathcal O$ is prime, so $\frac{n+\omega}p$ or its conjugate is in $\mathcal O=\Z+\Z\omega$, which is false since
the coefficient of $\omega$ is $\frac1p$. So $n^2+n+41$ is prime for $0\le n\le39$, Euler's famous observation
(Rabinowitsch proved in 1913 that this property is equivalent to $h_K=1$).
\end{solution}

\begin{solution}{ex:sqrt10}
\idea Primes above $2$ and $3$; congruences modulo $5$ show they are not principal; an element of norm $6$ links them.
\writeup $d_K=40$, $(r_1,r_2)=(2,0)$, $M_K=\frac2{4}\sqrt{40}=\sqrt{10}\approx3.16$. $2\mathcal O=P^2$ with $P=(2,\sqrt{10})$;
$3$ splits since $10\equiv1$ is a square modulo $3$: $3=Q\bar Q$. Modulo $5$, $x^2-10y^2\equiv x^2\in\{0,\pm1\}$, never
$\pm2$ or $\pm3$: so nothing has norm $\pm2$ or $\pm3$, and $P$, $Q$, $\bar Q$ are not principal. $N(4+\sqrt{10})=6$, so
$(4+\sqrt{10})=PQ$ after renaming, and $[Q]=[P]^{-1}=[P]$ (as $P^2=(2)$). Similarly $[\bar Q]=[Q]^{-1}=[P]$. So
$\operatorname{Cl}(K)=\{1,[P]\}$ has order $2$.
\end{solution}

\begin{solution}{ex:h10}
\idea Below the Minkowski bound lie only the primes above $2$ and $3$. Find their factorisations, then hunt for elements
of norm $2$ and $3$: if they exist, the primes are principal.
\writeup $(r_1,r_2)=(1,1)$, $d_K=-300$, $M_K=\frac4\pi\cdot\frac6{27}\sqrt{300}\approx4.90$: the class group is generated
by the primes above $2$ and $3$. At $2$: the index of $\Z[\theta]$ is $3$, prime to $2$, and $X^3-10\equiv X^3\pmod2$, so
$2\mathcal O=\mathfrak p_2^3$ with $\mathfrak p_2=(2,\theta)$ (Theorem~\ref{thm:kd}). At $3$ the index of $\Z[\theta]$ is
divisible by $3$, so we use $\beta$: its minimal polynomial $X^3-X^2-3X-3$ has discriminant $-300=d_K$, so
$\Z[\beta]=\mathcal O$, and modulo $3$ it is $X^2(X-1)$. So $3\mathcal O=\mathfrak q^2\mathfrak q'$ with $\mathfrak
q=(3,\beta)$ and $\mathfrak q'=(3,\beta-1)$, both of norm $3$.

Since the minimal polynomial of $\theta-a$ is $(X+a)^3-10$, $N(\theta-a)=10-a^3$; for $a=2$, $N(\theta-2)=2$. A principal
ideal of norm $2$ is a prime above $2$, so $(\theta-2)=\mathfrak p_2$. The constant term of the minimal polynomial of
$\beta$ gives $N(\beta)=3$, and $\beta\in\mathfrak q$, so $(\beta)=\mathfrak q$. Finally $\mathfrak q'=3\mathcal O\cdot
\mathfrak q^{-2}=(3/\beta^2)$ is principal too. All generators of $\operatorname{Cl}(K)$ are trivial: $h_K=1$.
\end{solution}

\begin{solution}{ex:discbound}
\idea Minkowski's inequality gives $|d_K|\ge b_n$; an induction on $n$ with a crude lower bound for $(1+\frac1n)^{2n}$
does the rest.
\writeup By Theorem~\ref{thm:dkbig} and $r_2\le\frac n2$, $|d_K|\ge b_n=(\frac\pi4)^{2r_2}(\frac{n^n}{n!})^2\ge(\frac\pi4)^n
(\frac{n^n}{n!})^2$. Then $b_2=\frac{\pi^2}{16}\cdot4=\frac{\pi^2}4=\frac\pi3\cdot\frac{3\pi}4$, and $\frac{b_{n+1}}{b_n}=
\frac\pi4\bigl(1+\frac1n\bigr)^{2n}\ge\frac\pi4\bigl(\frac32\bigr)^4=\frac{81\pi}{64}\ge\frac{3\pi}4$ for $n\ge2$ (the
sequence $(1+\frac1n)^n$ increases). By induction $b_n\ge\frac\pi3(\frac{3\pi}4)^{n-1}$. For $n=2$: $|d_K|\ge\frac{\pi^2}4
\approx2.47$, so $|d_K|\ge3$. For $n=3$ use the sharper $r_2\le1$: $|d_K|\ge(\frac\pi4)^2(\frac{27}6)^2\approx12.5$, so
$|d_K|\ge13$. For $n=4$, $|d_K|\ge(\frac\pi4)^4(\frac{256}{24})^2\approx43.3$, and for $n=5$, $\ge(\frac\pi4)^4(\frac{3125}
{120})^2\approx258$. So a field with $|d_K|<100$ has degree at most $4$. (The true minima are $3$, $23$, $117$ and $1609$
for degrees $2$ to $5$; in degree $4$, $117$ already exceeds $100$, a fact the bound cannot see.)
\end{solution}

\begin{solution}{ex:quadmin}
\idea The region $Q\le r$ is an ellipse; its area is computed by a linear change of variables, and Minkowski does the rest.
\writeup Write $Q(x,y)=a(x+\frac b{2a}y)^2+\frac{|D|}{4a}y^2$ ($a>0$ since $Q$ is positive definite). The linear map
$(x,y)\mapsto(\sqrt a(x+\frac b{2a}y),\frac{\sqrt{|D|}}{2\sqrt a}y)$ has determinant $\frac{\sqrt{|D|}}2$ and sends
$\{Q\le r\}$ onto the disc of radius $\sqrt r$, of area $\pi r$. So $\lambda(\{Q\le r\})=\frac{2\pi r}{\sqrt{|D|}}$. It is
compact, convex, symmetric; it has area $\ge4$ when $r=\frac2\pi\sqrt{|D|}$, so it contains $(x,y)\in\Z^2\setminus\{0\}$
with $0<Q(x,y)\le\frac2\pi\sqrt{|D|}$. Gauss's reduced forms have $a\le\sqrt{|D|/3}\approx0.577\sqrt{|D|}$, better than
$\frac2\pi\approx0.637$: for two variables the geometry of numbers is not optimal, but it works in every dimension.
\end{solution}

\begin{solution}{ex:simult}
\idea A thin, long box: long in the direction of $q$, thin in the directions of the errors; its volume is exactly $2^{k+1}$.
\writeup Let $S=\{x\in\R^{k+1}:|x_0|\le N^k,\ |x_0\alpha_i-x_i|\le\frac1N\ (1\le i\le k)\}$. It is compact, convex and
symmetric; it is the image of the box $[-N^k,N^k]\times[-\frac1N,\frac1N]^k$ under the linear map $(x_0,e)\mapsto(x_0,
x_0\alpha-e)$ of determinant $\pm1$, so $\lambda(S)=2N^k\cdot(\frac2N)^k=2^{k+1}$. Theorem~\ref{thm:minkowski}(b) gives
$(q,p_1,\dots,p_k)\in\Z^{k+1}\setminus\{0\}$ in $S$. For $N=1$ the statement is trivial ($q=1$, $p_i$ the integer part
of $\alpha_i$). For $N\ge2$, if $q=0$ then $|p_i|\le\frac1N<1$ forces $p_i=0$ for all $i$, which is excluded; so $q\ne0$. Replacing the point
by its opposite if needed, $1\le q\le N^k$ and $|q\alpha_i-p_i|\le\frac1N$.
\end{solution}

\section*{Sources and further reading}

\begin{itemize}
\item H.~Minkowski, \emph{Geometrie der Zahlen}, Teubner, Leipzig, 1896 (and 1910). The origin of the chapter.
\item D.~A.~Marcus, \emph{Number Fields}, 2nd ed., Springer, 2018, Chapter~5. The Minkowski bound, with many class
number computations; the source of Exercise~\ref{ex:discbound}.
\item J.~Neukirch, \emph{Algebraic Number Theory}, Springer, 1999, Chapter~I, \S\S4--6 and \S2 of Chapter~III (Hermite).
\item J.~W.~S.~Cassels, \emph{An Introduction to the Geometry of Numbers}, Springer Classics, 1997. The theory beyond
convex bodies: successive minima, Mahler's compactness theorem.
\item H.~Cohen, \emph{A Course in Computational Algebraic Number Theory}, GTM~138, Springer, 1993, \S5.4 and \S6.5.
Computing class groups in practice; tables of class numbers.
\item M.~Bhargava, ``The density of discriminants of quartic rings and fields'', \emph{Ann.\ of Math.}\ 162 (2005),
1031--1063. Counting fields by discriminant, the subject Hermite's theorem opens (Volume~14).
\item A.~K.~Lenstra, H.~W.~Lenstra, Jr.\ and L.~Lov\'asz, ``Factoring polynomials with rational coefficients'',
\emph{Math.\ Ann.}\ 261 (1982), 515--534. The LLL algorithm.
\item \emph{History.} E.~Bayer-Fluckiger, ``Hermann Minkowski, grand prix de l'Acad\'emie des sciences \`a 18 ans'',
Soci\'et\'e Math\'ematique de France, 2006: the 1881 question, the joint award of 2 April 1883, Smith's death, and Jordan's
letter. The English reaction: J.~J.~O'Connor and E.~F.~Robertson, ``Hermann Minkowski'' and ``Henry Smith'', MacTutor
History of Mathematics archive.
\end{itemize}
